\documentclass[11pt,a4paper]{scrartcl}
\usepackage[margin=1in,top=1.1in]{geometry}
\usepackage[T1]{fontenc}
\usepackage[utf8]{inputenc}
\usepackage{lmodern}
\usepackage{xcolor}
\usepackage{listings}
\usepackage{booktabs}
\usepackage{tabularx}
\usepackage{longtable}
\usepackage{array}
\usepackage{colortbl}
\usepackage{amsmath}
\usepackage{tikz}
\usepackage{tcolorbox}\tcbuselibrary{listings}
\usetikzlibrary{positioning,arrows.meta,shapes.geometric}
\usepackage{scrlayer-scrpage}
\usepackage[colorlinks=true,linkcolor=mainblue,urlcolor=accentblue,
            citecolor=mainblue,pdfborder={0 0 0}]{hyperref}

% the palette and boxes of the FOCEX agent manual
\definecolor{mainblue}   {RGB}{  0,  74, 153}
\definecolor{rungray}    {RGB}{110, 110, 110}
\definecolor{youteal}    {RGB}{  0, 128, 110}
\definecolor{accentblue} {RGB}{ 30, 120, 210}
\definecolor{lightblue}  {RGB}{210, 228, 255}
\definecolor{codebg}     {RGB}{245, 248, 252}
\definecolor{codeframe}  {RGB}{ 30, 120, 210}
\definecolor{notebg}     {RGB}{255, 250, 225}
\definecolor{notebar}    {RGB}{200, 130,   0}
\definecolor{dftgreen}   {RGB}{226, 242, 226}
\definecolor{fxorange}   {RGB}{255, 236, 214}
\definecolor{cmtgreen}   {RGB}{ 40, 130,  40}

\setkomafont{section}      {\color{mainblue}\large\bfseries}
\setkomafont{subsection}   {\color{mainblue}\normalsize\bfseries}
\setkomafont{descriptionlabel}{\color{mainblue}\bfseries}
\pagestyle{scrheadings}
\clearpairofpagestyles
\ohead{\small\textcolor{mainblue!60!black}{ALADYN GUI --- User's Guide}}
\ihead{\small\textcolor{mainblue!60!black}{\headmark}}
\cfoot{\pagemark}
\addtokomafont{pageheadfoot}{\normalfont\small}

\lstset{basicstyle=\ttfamily\small, breaklines=true, breakatwhitespace=true,
  backgroundcolor=\color{codebg}, frame=leftline, framerule=3pt, rulecolor=\color{codeframe},
  framesep=8pt, xleftmargin=14pt, commentstyle=\color{cmtgreen}\itshape, upquote=true,
  columns=fullflexible, keepspaces=true, aboveskip=8pt, belowskip=4pt, morecomment=[l]{\#}}
\newlength{\boxwidth}
\newenvironment{notebox}{%
  \par\vspace{6pt}\setlength{\tabcolsep}{5pt}%
  \setlength{\boxwidth}{\dimexpr\linewidth - 4pt - 2\tabcolsep}%
  \noindent\begin{tabular}{!{\color{notebar}\vrule width 4pt}>{\columncolor{notebg}}p{\boxwidth}}%
  \rule{0pt}{11pt}\textbf{\textcolor{notebar}{Note.}}\enspace\ignorespaces}{%
  \rule[-4pt]{0pt}{4pt}\end{tabular}\par\vspace{6pt}}
\setlength{\emergencystretch}{3em}
\newtcblisting{doit}[1][]{listing only, colback=lightblue!35, colframe=mainblue,
  boxrule=0.8pt, arc=2pt, left=6pt, right=4pt, top=3pt, bottom=3pt, fonttitle=\bfseries\small,
  title={Do this#1}, listing options={basicstyle=\ttfamily\small\color{mainblue}, breaklines=true,
  breakatwhitespace=true, columns=fullflexible, keepspaces=true, upquote=true, frame=none,
  backgroundcolor={}, xleftmargin=0pt, framesep=0pt, aboveskip=0pt, belowskip=0pt,
  commentstyle=\color{mainblue!55}\itshape}}
\newcommand{\thead}{\rowcolor{lightblue}}
\newcommand{\cmd}[1]{\texttt{\detokenize{#1}}}
\newcommand{\btn}[1]{\fcolorbox{mainblue}{lightblue!40}{\small\textsf{#1}}}

\title{%
  \vspace*{-10pt}{\color{mainblue}\rule{\linewidth}{2pt}}\\[8pt]%
  \textbf{\Huge ALADYN GUI}\\[4pt]\large User's Guide\\[6pt]%
  \normalsize From a POSCAR to phonons and the thermal conductivity, in a browser:
  VASP $\cdot$ SC-SNAPS $\cdot$ FOCEX $\cdot$ THERMACOND\\[6pt]{\color{mainblue}\rule{\linewidth}{0.5pt}}}
\author{Keivan Esfarjani}
\date{v0.1 $\cdot$ November 2026\\[2pt]\small\url{https://github.com/KeivanS/Anharmonic-lattice-dynamics} (\texttt{GUI})}

\begin{document}
\maketitle
\tableofcontents
\newpage

% ============================================================
\section{What the GUI does}
% ============================================================
The ALADYN GUI (\cmd{aladyn-gui}) takes you, in a web browser, through a whole lattice-dynamics
calculation: relaxation of the crystal, SCF and Born charges, supercell snapshots (SC-SNAPS), one DFT
run per snapshot, the force-constant fit (FOCEX) with its phonon plots, and the lattice thermal
conductivity (THERMACOND). You upload a POSCAR and press buttons; a panel on the left always says which
step comes next and what to do.

Underneath it is the FOCEX agent (\cmd{FOCEX/AGENT/focex-agent.py}, its own manual
\cmd{focexagent_manual.pdf}): every GUI project is an agent project folder, the GUI writes the agent's
instruction file and runs the agent's tasks one at a time. What the GUI adds:
\begin{itemize}
\item a choice, for each DFT stage and for THERMACOND, of \emph{where} it runs: on this computer, as a
  bundle for a SLURM cluster, or submitted directly when the GUI itself runs on the cluster
  (Section~\ref{sec:where});
\item editors for the input files, with a comment after every line and the values that are not known yet
  marked \cmd{PLACEHOLDER} (Section~\ref{sec:params});
\item stops before FOCEX and before THERMACOND, so that you review their input files first;
\item the plots and the SC-SNAPS displacement histograms in the page.
\end{itemize}

\begin{figure}[h]
\centering
\begin{tikzpicture}[node distance=2.5mm and 2.2mm, every node/.style={font=\footnotesize},
  s/.style={draw, rounded corners=2pt, align=center, minimum height=10mm, text width=17mm},
  d/.style={s, fill=dftgreen}, f/.style={s, fill=fxorange}, l/.style={s, fill=lightblue!50},
  a/.style={-{Stealth[length=2mm]}, thick, mainblue}]
  \node[l] (p) {Project\\POSCAR};
  \node[d, right=of p] (r) {Relaxation};
  \node[d, right=of r] (s) {SCF\\+ DFPT};
  \node[l, right=of s] (n) {Snapshots\\(here)};
  \node[d, right=of n] (v) {Supercell\\DFT};
  \node[f, right=of v] (x) {FOCEX\\+ plots};
  \node[f, right=of x] (t) {THERMA-\\COND};
  \foreach \i/\j in {p/r, r/s, s/n, n/v, v/x, x/t} \draw[a] (\i) -- (\j);
  \node[below=4mm of s, align=center, text width=70mm] (w)
    {\textbf{green and THERMACOND:} run here $\cdot$ bundle for the cluster $\cdot$ submit here};
  \draw[densely dotted, mainblue] (r.south) -- (w.north west);
  \draw[densely dotted, mainblue] (v.south) -- (w.north east);
  \draw[densely dotted, mainblue] (t.south) |- (w.east);
\end{tikzpicture}
\caption{The seven steps of the guide panel. The GUI stops before FOCEX and THERMACOND (orange) until you
start them; blue steps are quick and always run where the GUI runs.}
\end{figure}

% ============================================================
\section{Installation and start}
% ============================================================
The GUI is part of the ALADYN repository (\cmd{GUI/aladyn-gui.py}). It needs what the FOCEX agent needs
(Python with numpy, spglib and pymatgen; VASP-Flow; SC-SNAPS and FOCEX built; a POTCAR library) and the
Python package Flask. Set these up once as described in the FOCEX agent manual (\cmd{focex-configure.sh}).

\begin{doit}[: once, and after every git pull]
cd ~/PROJECTS/Anharmonic-lattice-dynamics
./configure               # this computer's settings (--install the first time); checks Flask too
make                         # also links aladyn-gui into ~/ALADYN/bin
pip install flask            # in the agent's Python environment, if it is not there
\end{doit}

\subsection*{On your computer}
\begin{doit}
source ~/myenv/bin/activate              # the agent's Python environment
cd ~/calc && aladyn-gui                  # projects go in this folder
\end{doit}
Then open \url{http://localhost:5070}. Options: \cmd{--workdir DIR} (folder of the projects, default the
current one), \cmd{--port N} (default 5070), \cmd{--host} (default 127.0.0.1: only this computer can
connect).

\subsection*{On a cluster}
Start it on the login node and reach it from your computer through an ssh tunnel. Then the DFT runs can be
submitted from the page, and nothing has to be copied back and forth.
\begin{doit}
ssh -L 5070:localhost:5070 you@cluster       # on your computer; keep this open
cd /scratch/$USER && aladyn-gui              # on the cluster
\end{doit}
Open \url{http://localhost:5070} on your computer. The button \btn{Submit here} is active when the GUI finds
\cmd{sbatch}.

% ============================================================
\section{The page}
% ============================================================
\textbf{The guide} (left, always visible) lists the steps in order, each marked \emph{done}, \emph{next},
\emph{waiting} (its inputs are ready and the runs are outstanding), \emph{failed} or \emph{later}; a step
done outside (a relaxed POSCAR, results of a cloned project) says \emph{provided}. Below it,
\textbf{What to do next} names the button to press. It is updated every few seconds while a job runs.

\medskip\noindent
\begin{tabularx}{\linewidth}{@{}>{\bfseries}p{0.18\linewidth}X@{}}
\thead Tab & What it holds \\ \midrule
Projects & open a project; a new project from a POSCAR; a new project from the open one \\
Run & the run buttons, the import of results, the bundles made so far, the cluster settings \\
Snapshots & the SC-SNAPS summary (supercell, temperature, displacements) and two histograms \\
Params & the five FOCEX input files \\
Results & the PDF plots of FOCEX (\cmd{05_focex}, also \cmd{NOBORN/} and \cmd{LASSO/}) and THERMACOND \\
THERMACOND & the \cmd{latdyn.params} of THERMACOND and its run buttons \\
Setup & the whole instruction file of the project (checked before it is saved) \\
Log & the output of the running job and \cmd{report.log} (one entry per task) \\
\bottomrule
\end{tabularx}

% ============================================================
\section{A calculation, step by step}
% ============================================================
\paragraph{1. New project.} On \emph{Projects}: a name, the POSCAR (paste it or choose the file), and the
main settings:

\medskip\noindent
\begin{tabularx}{\linewidth}{@{}>{\bfseries}p{0.27\linewidth}X@{}}
\thead Field & Meaning (instruction keyword) \\ \midrule
Functional & PBE, PBEsol, LDA, R2SCAN or HSE06 (\cmd{FUNCTIONAL}) \\
Accuracy & Accurate: ENCUT 600 eV, EDIFF $10^{-7}$; Normal (\cmd{ACCURACY}) \\
DFPT & Born charges: auto (when the SCF shows a gap), on, off (\cmd{DFPT}) \\
DFT code & VASP (others later) \\
Force-constant ranks & harmonic; harmonic + cubic (needed by THERMACOND); + quartic (\cmd{RANKS}) \\
Shells for ranks 2..8 & neighbour shells of each rank (\cmd{RANGE}) \\
Supercell & SMALL, LARGE or 9 integers (\cmd{SUPERCELL}) \\
Snapshots & number of snapshots (\cmd{SNAPSHOTS}) \\
Temperature & of the snapshots; empty: 30\,K tuned to an rms displacement of 0.02\,\AA \\
THERMACOND job & 0 none, 1 linewidths, 2 and 3 thermal conductivity (\cmd{THERMACOND_JOB}) \\
MPI ranks here & VASP and THERMACOND ranks on this computer \\
More instruction lines & any keyword of the agent or VASP-Flow (\cmd{MAGMOM}, \cmd{GGA_U}, INCAR blocks) \\
\bottomrule
\end{tabularx}

\medskip
Tick \emph{This POSCAR is already relaxed} only if it was relaxed with these same DFT settings: the
supercell runs must use the settings of the relaxation. Otherwise the GUI relaxes it first.
\btn{Create project} makes the folder, the POTCAR, the relaxation inputs, and the five draft input files.

\paragraph{2. Relaxation.} On \emph{Run}: \btn{Run here} (a small cell relaxes in minutes on a
workstation), \btn{Make bundle} or \btn{Submit here}. A relaxed POSCAR: \btn{Skip: POSCAR is relaxed}.
Then \btn{Continue}: the relaxed cell becomes \cmd{POSCAR0}, the cell of every later step.

\paragraph{3. SCF and DFPT.} The same buttons. With DFPT both go into one bundle or submission; DFPT starts
after the SCF and reads its charge density.

\paragraph{4. Snapshots.} \btn{Continue} makes the snapshots (SC-SNAPS) and one VASP folder per snapshot,
here; see the \emph{Snapshots} tab. It stops with the supercell runs waiting.

\paragraph{5. Supercell DFT.} The expensive step: usually \btn{Make bundle} (or \btn{Submit here} on the
cluster). All snapshots are submitted at once; a last job writes \cmd{FORCEDISP1}, \cmd{POSCAR1} and
\cmd{stress.params}. Import the results, then \btn{Continue}: the guide moves to FOCEX.

\paragraph{6. FOCEX.} On \emph{Params}, read the five files (Section~\ref{sec:params}), change what you
want, then \btn{Save and run FOCEX}. FOCEX runs here (minutes) and the plots follow (\emph{Results}).

\paragraph{7. THERMACOND.} On \emph{THERMACOND}, check line 8 (the job), line 1 (q mesh) and line 6
(temperatures), then \btn{Save and run here}, \btn{Save and make bundle} or \btn{Save and submit here}. The
thermal conductivity is in \cmd{07_thermacond/kappa_vs_T.dat}; the table of \cmd{report.log} (\emph{Log})
gives it at a few temperatures.

% ============================================================
\section{Where the runs go: here, bundle, submit}\label{sec:where}
% ============================================================
\begin{description}
\item[Continue] does the quick steps here (POSCAR0, snapshots, folders, reading results that came back)
  and stops at the next DFT stage with its inputs ready (the guide then says \emph{waiting}). It never
  starts VASP.
\item[Run here] runs the next DFT stage on this computer, with the \emph{MPI ranks here}.
\item[Make bundle] packs the stage as \cmd{bundles/<project>_<task>_<date>.tar.gz} (link on the
  \emph{Run} tab) for a SLURM cluster.
\item[Submit here] submits the same scripts with \cmd{sbatch} (the GUI runs on the cluster). When the jobs
  are done, \btn{Continue}.
\end{description}

\subsection{On the cluster, with a bundle}
\begin{doit}
tar xzf Cu2O_vasp_snaps_20261107-214157.tar.gz    # where the project should live (scratch)
cd Cu2O && ./submit_all.sh
squeue -u $USER                                   # until the jobs are gone
./pack_results.sh                                 # -> results_Cu2O_<date>.tar.gz
\end{doit}
Copy \cmd{results_*.tar.gz} back, choose it under \emph{Import results} and press \btn{Import results},
then \btn{Continue}. The same folder name is not needed: the top folder of the archive is dropped.

\medskip\noindent
\begin{tabularx}{\linewidth}{@{}>{\ttfamily}p{0.27\linewidth}X@{}}
\thead \normalfont\bfseries In the bundle & \textbf{What it is} \\ \midrule
submit\_all.sh & submits every waiting folder; DFPT waits for the SCF; the supercells all at once,
  then the collect job \\
<folder>/job.sh & the SLURM job of that folder (header from the cluster settings) \\
<folder>/INCAR, KPOINTS, POSCAR, POTCAR & the VASP inputs (the POTCAR is built on your computer) \\
04\_supercells/collect.py & writes FORCEDISP1, POSCAR1, stress.params and collect.log \\
04\_supercells/collect\_job.sh & runs collect.py after all supercell jobs (30 min, one task) \\
06\_dfpt/gap\_check.py & with DFPT auto: the job ends at once when the SCF shows no gap \\
pack\_results.sh & packs the files to bring back \\
README\_HPC.txt & these commands \\
\bottomrule
\end{tabularx}

\medskip
\cmd{pack_results.sh} packs, for the relaxation, SCF and DFPT, the OUTCAR, CONTCAR, OSZICAR, EIGENVAL,
\cmd{vasp.out} and the Born-charge files; for the supercells only \cmd{FORCEDISP1}, \cmd{POSCAR1},
\cmd{stress.params}, \cmd{collect.log} and each run's OSZICAR and \cmd{vasp.out}
(\cmd{./pack_results.sh --outcars} adds the OUTCARs); for THERMACOND every file of
\cmd{07_thermacond} but its scratch folder \cmd{V3DELTA}.

\begin{notebox}
\cmd{collect.py} does what \cmd{read_outcar.x} does, in Python with the standard library, so nothing has
to be compiled on the cluster. It checks every run (finished, converged), the net force and the positions,
and writes nothing if a run is not usable (\cmd{collect.log} names it). On the 21 Cu$_2$O snapshots of
2026-11-07 it gave the same numbers as \cmd{read_outcar.x}.
\end{notebox}

\subsection{Cluster settings}
On the \emph{Run} tab, kept in the project's \cmd{hpc.json}. A new project starts from the VASP-Flow profile
\cmd{FOCEX_PROFILE} that \cmd{./configure} drafted on the cluster (partition, account, nodes, tasks, time,
modules, VASP command), else from the defaults below:

\medskip\noindent
\begin{tabularx}{\linewidth}{@{}>{\bfseries}p{0.32\linewidth}X@{}}
\thead Setting & Meaning (default) \\ \midrule
Partition, Account & \texttt{\#SBATCH -{}-partition}, \texttt{-{}-account} (standard, none) \\
Nodes, Tasks per node & per VASP job (1, 40); also the ranks used for KPAR \\
Walltime (VASP), (THERMACOND) & \texttt{\#SBATCH -{}-time} (12:00:00, 24:00:00) \\
Memory & \texttt{\#SBATCH -{}-mem}, optional \\
MPI launcher & \cmd{srun}, or e.g.\ \cmd{mpirun} (then \cmd{-np} nodes $\times$ tasks) \\
VASP / THERMACOND binary & on the cluster (\cmd{vasp_std}, \cmd{~/ALADYN/bin/tcond-2.2.x}) \\
python3 on the cluster & for collect.py, gap\_check.py, set\_efield.py, extract\_born.py \\
Lines before VASP & e.g.\ \cmd{module load vasp/6.4.2} \\
More \#SBATCH lines & anything else your cluster wants \\
\bottomrule
\end{tabularx}

\medskip
KPAR and NCORE in the INCARs are adjusted to the number of ranks the job runs with (KPAR divides the ranks).

% ============================================================
\section{The input files}\label{sec:params}
% ============================================================
The \emph{Params} tab shows \cmd{structure.params}, \cmd{latdyn.params}, \cmd{kpbs.params},
\cmd{dielectric.params} and \cmd{default.params} as FOCEX will get them, each line followed by a comment
that says what it is. Lines marked \cmd{PLACEHOLDER} are filled in by a later step or need your values:
\begin{itemize}
\item \cmd{structure.params}, before the relaxation: lattice and positions of the unrelaxed POSCAR. When
  the relaxation is done, they are replaced by those of \cmd{POSCAR0}; your other lines (ranks,
  invariances, shells, \dots) are kept.
\item \cmd{dielectric.params}, before DFPT: zeros and \cmd{born_flag} 0. After DFPT the Born charges are
  there. You may also type $\varepsilon_\infty$ and $Z^*$ yourself (e.g.\ from a database) and set line 1
  to 1.
\end{itemize}
A file you change is kept in the instruction file and used as it is; \btn{Back to the generated file}
undoes your changes (save afterwards). A saved \cmd{dielectric.params} is used instead of the DFPT result.

The \emph{THERMACOND} tab holds the \cmd{latdyn.params} of \cmd{07_thermacond}: before FOCEX has run it
is FOCEX's file with the THERMACOND lines of your settings; afterwards the one FOCEX used. The lines that
matter most: 1 the q mesh, 6 the temperatures (Tmin, Tmax, number), 8 the job, 9 the solver of job 2,
11 isotope scattering, 13 the sample length for boundary scattering.

% ============================================================
\section{More projects from one}
% ============================================================
\emph{New project from the open one} (\emph{Projects} tab) keeps the relaxed cell, the SCF and the DFPT
(they are marked \emph{provided}) and makes the snapshots again with another temperature, supercell or
number of snapshots. A second supercell (\cmd{FORCEDISP2} for one FOCEX fit) is a separate project.

% ============================================================
\section{The project folder}
% ============================================================
The folder is that of the FOCEX agent (its manual, Section ``The project folder'':
\cmd{01_relax}, \cmd{02_scf}, \cmd{03_snaps}, \cmd{04_supercells}, \cmd{05_focex}, \cmd{06_dfpt},
\cmd{07_thermacond}, \cmd{report.log}, \cmd{state.json}). The GUI adds:

\medskip\noindent
\begin{tabularx}{\linewidth}{@{}>{\ttfamily}p{0.3\linewidth}X@{}}
\thead \normalfont\bfseries File & \textbf{What it is} \\ \midrule
gui.json & the MPI ranks here; when the input files were saved \\
hpc.json & the cluster settings \\
bundles/ & the bundles made, and the imported results (\cmd{imported_*}) \\
submit\_all.sh, pack\_results.sh, README\_HPC.txt & of the last stage made ready \\
<folder>/job.sh & SLURM job of a waiting folder \\
\bottomrule
\end{tabularx}

% ============================================================
\section{When something goes wrong}
% ============================================================
\begin{description}
\item[A step is marked failed] The guide quotes the message of \cmd{report.log}. Fix the cause (the
  instruction file on \emph{Setup}, or the INCAR of that folder) and press \btn{Continue}. The agent's
  \cmd{diagnose.py} explains most VASP failures and writes a fix script (its command is in the guide).
\item[The step stays \emph{waiting}] The results are not in the folders yet: import them, or (GUI on the
  cluster) wait until \cmd{squeue} is empty, then \btn{Continue}.
\item[collect.py wrote no FORCEDISP1] \cmd{collect.log} names the runs that did not finish or converge.
  Rerun them on the cluster (\cmd{cd 04_supercells/snap_NNN && sbatch job.sh}), then
  \cmd{python3 collect.py} in \cmd{04_supercells} and \cmd{./pack_results.sh}.
\item[No POTCAR] The POTCAR library was not found (\cmd{VASP_POTCAR_DIR}); put a POTCAR in the project
  folder.
\item[Born charges not used] The DFPT result failed the agent's checks (charges not summing to zero,
  $\varepsilon_\infty$ outside 1--100, VASP's Zener warning); \cmd{report.log} says why. Type the values
  in \cmd{dielectric.params}, or redo DFPT.
\end{description}

% ============================================================
\section{Version history}
% ============================================================
\begin{description}
\item[0.1] 2026-11-08: first version, with the FOCEX agent 1.6 (\cmd{VASP_RUN} and
  \cmd{THERMACOND_RUN} external) and \cmd{FOCEX/AGENT/bundle.py}. Tested on Cu$_2$O through all steps with
  the GUI's own requests (no browser, no cluster yet).
\end{description}
\end{document}
